% fig:wt-net --- the Petri net the elaborated ACFG lowers to, and the soundness
% gate. Faithful to sec:wt-net of chapters/07b-compilation-walkthrough.tex:
%   places  = array regions a worker owns, transfer channels, barriers;
%   transitions = host load_image, 4 img_in sends + 4 matching receives,
%                 each worker's blur3 firings, the two barriers,
%                 4 img_out sends + receives, host save_image.
%   channel capacities: 2 for the async img_in channels, 1 for the sync img_out.
% The four compute workers are a 4-fold-symmetric column; drawn once inside a
% plate ("x4") for legibility, with the host spine across the top. The two
% barriers are 5-way (host + w0..w3).
\begin{figure}[tbp]
  \centering
  \resizebox{\textwidth}{!}{%
  \begin{tikzpicture}[
      font=\small,
      place/.style={draw, circle, minimum size=9mm, inner sep=1pt, align=center,
                    fill=green!10},
      chan/.style={draw, circle, minimum size=10mm, inner sep=1pt, align=center,
                   fill=blue!10, very thick},
      tr/.style={draw, fill=gray!75, minimum width=2.2mm, minimum height=10mm,
                 inner sep=0pt},
      bar/.style={draw, fill=orange!70, minimum width=2.6mm, minimum height=34mm,
                  inner sep=0pt},
      a/.style={-{Latex[length=1.8mm]}},
      lbl/.style={font=\scriptsize, align=center},
    ]
    % ---- host spine (y = 2.6) ----
    \def\hy{2.6}
    \node[tr] (tload) at (0,\hy) {};   \node[lbl, above=0.5mm of tload] {\texttt{load\_image}};
    \node[place] (pin) at (1.5,\hy) {\texttt{img\_in}\\\scriptsize @host};
    \node[tr] (tsend) at (3.0,\hy) {}; \node[lbl, above=0.5mm of tsend] {send};
    \node[tr] (trout) at (12.6,\hy) {};\node[lbl, below=0.5mm of trout] {recv};
    \node[place] (pout) at (14.1,\hy) {\texttt{img\_out}\\\scriptsize @host};
    \node[tr] (tsave) at (15.6,\hy) {};\node[lbl, above=0.5mm of tsave] {\texttt{save\_image}};

    % ---- a single worker column (y = 0), enclosed by a plate ----
    \node[chan] (cin) at (4.5,1.3) {cap\,2};
    \node[lbl, left=0.5mm of cin] {\texttt{img\_in}\\chan};
    \node[tr] (trin) at (6.0,0) {};   \node[lbl, below=0.5mm of trin] {recv};
    \node[place] (pband) at (7.4,0) {band\\\scriptsize $+$halo};
    \node[tr] (tblur) at (8.8,0) {}; \node[lbl, below=0.5mm of tblur] {\texttt{blur3}};
    \node[place] (pob) at (10.0,0) {out\\\scriptsize band};
    \node[tr] (tso) at (11.1,0) {};  \node[lbl, below=0.5mm of tso] {send};
    \node[chan] (cout) at (12.6,1.3) {cap\,1};
    \node[lbl, right=0.5mm of cout] {\texttt{img\_out}\\chan};

    % plate around the worker column
    \begin{scope}[on background layer]
      \node[draw, dashed, rounded corners, fill=yellow!8, inner sep=4mm,
            fit=(cin)(trin)(pband)(tblur)(pob)(tso)(cout),
            label={[font=\footnotesize\bfseries, fill=yellow!18, rounded corners,
                    inner sep=2pt, yshift=-2mm]below:$\times\,4$ --- one column per
                    worker (\texttt{w0}--\texttt{w3}): bands $1..5$ / $5..9$ /
                    $9..12$ / $12..15$}]
            (plate) {};
    \end{scope}

    % ---- the two 5-way barriers ----
    \node[bar] (b1) at (5.0,1.4) {};
    \node[lbl, orange!60!black] at (5.0,3.7) {barrier 1\\\scriptsize (pre-compute, 5-way)};
    \node[bar] (b2) at (11.9,1.4) {};
    \node[lbl, orange!60!black] at (11.9,3.7) {barrier 2\\\scriptsize (post-compute, 5-way)};

    % ---- arcs ----
    \draw[a] (tload) -- (pin);
    \draw[a] (pin) -- (tsend);
    \draw[a] (tsend) -- (cin);            % host pushes into the img_in channel
    \draw[a] (cin) -- (trin);             % worker receives
    \draw[a] (trin) -- (pband);
    \draw[a] (pband) -- (tblur);
    \draw[a] (tblur) -- (pob);
    \draw[a] (pob) -- (tso);
    \draw[a] (tso) -- (cout);             % worker pushes into img_out channel
    \draw[a] (cout) -- (trout);           % host receives
    \draw[a] (trout) -- (pout);
    \draw[a] (pout) -- (tsave);
    % barrier synchronisation arcs (host + worker both meet each barrier)
    \draw[a, orange!60!black] (tsend) -- (b1);
    \draw[a, orange!60!black] (b1) -- (trin);
    \draw[a, orange!60!black] (tso) -- (b2);
    \draw[a, orange!60!black] (b2) -- (trout);
    % host participates in BOTH barriers: it passes barrier 1, idles, then meets
    % barrier 2 (the back-to-back "barrier; barrier" in the host's event list).
    \draw[a, orange!60!black] (b1.north) to[out=70,in=110] (b2.north);
  \end{tikzpicture}}
  \caption[The stencil's Petri net and soundness gate]{The Petri net the
    elaborated ACFG lowers to. \emph{Places} (circles) model state --- the array
    regions a worker owns and the transfer channels --- and \emph{transitions}
    (bars) model events: the host's \texttt{load\_image}, the four
    \texttt{img\_in} sends and matching receives, each worker's \texttt{blur3}
    firings, the four \texttt{img\_out} sends and receives, and the host's
    \texttt{save\_image}. The four compute workers form a 4-fold-symmetric
    column, drawn once inside the dashed plate. Channel places (marked
    ``cap'') carry the capacity their buffer policy sets: \textbf{2} for the
    asynchronous \texttt{img\_in} channels, \textbf{1} for the synchronous
    \texttt{img\_out} channels. The two orange bars are the 5-way (host $+$
    \texttt{w0}--\texttt{w3}) pre- and post-compute barriers; the orange arcs
    are their barrier-place arcs (no participant proceeds until all five have
    arrived), and the arc arching over the top is the host idling from the
    first barrier to the second. The gate checks, in order, conflict-freedom
    (no place offers a token to two competing consumers), then --- on one
    replayed firing order --- boundedness (no \texttt{img\_in} channel ever
    exceeds 2) and deadlock-freedom, before any code is emitted.}
  \label{fig:wt-net}
\end{figure}
